\section{Forces as Twists}

Matter is knots in the web. Forces are what move those knots, push and pull, attract and repel. In this theory a force is, at bottom, a twist that a vibe's state picks up as it crosses a note. This chapter gives the gist of how the familiar forces come out, without the heavy machinery.

\subsection{A force is a twist along a note}

When a pattern spreads from one vibe to its neighbor across a note, the state can get a little twist on the way, a small adjustment in how it is oriented. If every note in a region twists things the same way, patterns crossing that region get systematically deflected. That systematic deflection is a force. This is the same idea physicists already use for electromagnetism and the nuclear forces, just expressed in the crystal: a force is a field of twists living on the notes, and particles feel the force by picking up those twists as they propagate.

So forces are not a new ingredient either. They are a texture in the fills and notes, a pattern of how crossings twist, and matter responds to that texture as it moves.

\subsection{Confinement: why quarks are never alone}

One striking thing the crystal reproduces is confinement, the reason you never find a lone quark. Quarks always come bound together, into protons and neutrons, and no amount of pulling frees a single one. Pull hard enough and the energy you pour in just makes new quarks, which pair up again. This is mysterious from the outside and famously hard to derive.

In the crystal, the twist-field between two quarks does not spread out and weaken with distance the way ordinary electric pull does. Instead it bunches into a tight tube, and a tube costs the same per length no matter how long, so pulling the quarks apart costs ever more energy without end, until it is cheaper to snap the tube and make a new pair. The simulation shows exactly this bunching and the resulting permanent imprisonment of the quarks. Confinement comes out, not put in.

\subsection{The forces fall into families by spin}

Recall that patterns carry spin, a locked-in twist. The carriers of force are patterns of a particular spin (spin one for light and the nuclear forces, spin two for gravity), and the matter they push is patterns of another (spin one-half). The crystal naturally hosts patterns of these different spins, so the cast of force-carriers and matter particles sorts itself by spin the way the real particle zoo does. Light, the carrier of electromagnetism, is one such pattern, a massless spin-one ripple, and it comes out with no mass exactly as it should.

\subsection{Where the force laws come from}

The deepest version of this story is that you do not put the force laws in by hand at all. You write down one quantity for the whole crystal, a kind of total tension of the configuration, and ask how it behaves under gentle ripples. The equations for light and even for gravity drop out as the gentle-ripple limit of that one quantity. So the separate-looking forces are facets of a single underlying tension on a single web, which is about as unified as a theory of forces can hope to be.

We will not work the machinery. The point for this walkthrough is the shape of the answer. Forces are twists on the notes, matter feels them by crossing those notes, confinement and masslessness come out correctly, and the force laws themselves descend from one quantity on the crystal rather than being separately assumed.

\textbf{Takeaway:} a force is a field of twists living on the notes, and matter feels a force by picking up those twists as it moves. Confinement of quarks and the masslessness of light both come out of the crystal, and the force laws descend from a single tension on the whole web rather than being put in by hand.
